\documentclass[10pt,a4paper]{amsart}
\usepackage[margin=19mm]{geometry}
\usepackage{amsmath,amssymb,mathtools}
\usepackage[scr=rsfs,cal=euler]{mathalfa}
\usepackage{xfrac}
\usepackage[unicode,hidelinks,bookmarks=false]{hyperref}
\newcommand{\ie}{i.e.\ }
\newcommand{\cf}{cf.\ }
\newcommand{\resp}{respectively\ }
\newcommand{\Subsection}{\subsection}
\providecommand{\nobreakdash}{\nobreak\hbox{-}}
\setlength{\emergencystretch}{2em}
\multlinegap0pt
\usepackage{amssymb}
\newtheorem{lem}{Lemma}
\newtheorem{thm}{Theorem}
\title{On Sharp Estimates of Derivatives of Even Order\\ in Sobolev Spaces}
\author{T.~A.~Garmanova and I.~A.~Sheipak}
\date{Source: arXiv:2606.22646v1 (CC BY 4.0)}
\begin{document}
\maketitle

\noindent\emph{Source and licence.} On Sharp Estimates of Derivatives of Even Order\\ in Sobolev Spaces. Version arXiv:2606.22646v1 (CC BY 4.0), distributed by arXiv under the Creative Commons Attribution 4.0 International licence, http://creativecommons.org/licenses/by/4.0/, as recorded in the arXiv licence metadata for this exact version and shown on the abstract page https://arxiv.org/abs/2606.22646v1. Full licence notice: \texttt{licenses/arxiv-2606.22646-CC-BY-4.0.txt}.\par\smallskip

\noindent\textit{Editorial scope.} Counted unit: the article's thm thm:derAPP (source0.tex lines 747--763). The article's complete original proof of it is delivered with it (source0.tex lines 765--839, one \texttt{proof} environment of 75 lines). No mathematics is written, repaired or re-derived: the statement and the proof are the article's own lines, in the article's own notation, and no retained line is substituted or re-flowed.  Retained uncounted as the source-local context the proof needs: lines 429--432; lines 485--488; lines 539--547.  Nothing was omitted from any retained span, so the package carries no omission record.  No retained line contains a citation, so the wrapper carries no bibliography and no bibliography entry.  No retained line contains a figure environment, an image-inclusion command or drawing code, so no image is delivered and the package declares no asset dependency.  Editorial formatting only, in addition: an AMS \texttt{amsart} preamble that restates the article's own declarations and macros used by the retained lines, namely the environments \texttt{lem}, \texttt{thm} on separate counters, together with the standard packages those lines need.  The retained environments print in this document as Lemma 1, Theorem 1 in order of appearance, whereas the article prints the same items with the numbering of its own full document.  The complete native source of the article as distributed in the arXiv e-print for this exact version is bundled unchanged as \texttt{sources/203368/source0.tex}.
\begin{equation}
\label{eq:Pdifur}
(x-x^2)y''+(2x-1)(m-1)y'+(n+m)(n-m+1)y=0.
\end{equation}

\begin{equation}
\label{eq:Prec}
2(2n+1) P_{n}^{(-m)}=P_{n+1}^{(-m+1)}-P_{n-1}^{(-m+1)}.
\end{equation}

\begin{lem}
\label{lem:APP}
 For the functions
$A^2_{n,k}(x)$,
the following relation holds{\rm:}
$$
\dfrac{d}{dx}(A^2_{n,k})=-P_{n-1}^{(k-n+1)}\cdot P_{n}^{(k-n+1)}.
$$
\end{lem}

\begin{thm}
\label{thm:derAPP}
 The values of the quantities
$A^2_{n,k}$
at their local maximum points
lying
on the closed interval
$[0,{1}/{2}]$
constitute a nondecreasing sequence
and those
at the local maximum points
lying
on the closed interval
$[\frac{1}{2},1]$
constitute a nonincreasing sequence
(see the figure).
\end{thm}

\begin{proof}
 For each fixed
$n$
and
$k$,
we consider the function
$$
B_{n,k}(x) = A^2_{n,k}(x) + (f(x)+g(x)) (P_{n}^{(k-n+1)}(x))^2
$$
with
some,
so far arbitrary, differentiable functions~$f$ and~$g$.
 The concrete form of these functions will be indicated in what follows.
 Then
\begin{align*}
%\begin{multline*}
B'_{n,k}(x) &= (A^2_{n,k})' + (f'(x)+g'(x)) (P_{n}^{(k-n+1)})^2 + 2 (f(x)+g(x))
P_{n}^{(k-n+1)} P_{n}^{(k-n+2)}\\
&=-P_{n-1}^{(k-n+1)} P_{n}^{(k-n+1)} + (f'(x)+g'(x)) (P_{n}^{(k-n+1)})^2 + 2 (f(x)+g(x))
P_{n}^{(k-n+1)} P_{n}^{(k-n+2)}\\
&=P_{n}^{(k-n+1)} \left(-P_{n-1}^{(k-n+1)} + g'(x) P_{n}^{(k-n+1)} + 2
(f(x)+g(x))P_{n}^{(k-n+2)} \right) + f'(x) (P_{n}^{(k-n+1)})^2.
%\end{multline*}
\end{align*}
 It follows from the recurrence relation~\eqref{eq:Prec}
and Lemma~\ref{lem:APP} that
$$
-P_{n-1}^{(k-n+1)} = 2(2n+1)P_{n}^{(k-n)}-P_{n+1}^{(k-n+1)} = 2(2n+1)P_{n}^{(k-n)}-2(k+1)
P_{n}^{(k-n)} - (2x-1) P_{n}^{(k-n+1)}.
$$
 Thus,
\begin{multline*}
B'_{n,k}(x) = P_{n}^{(k-n+1)} \left(2(2n-k)P_{n}^{(k-n)}+ g'(x) P_{n}^{(k-n+1)} + 2
(f(x)+g(x))P_{n}^{(k-n+2)} \right) +\\+
f'(x) (P_{n}^{(k-n+1)})^2+(1-2x)(P_{n}^{(k-n+1)})^2.
\end{multline*}
 We choose
$f(x)$
and
$g(x)$
so that the first summand
is equal to zero.
 Taking into account the differential equation~\eqref{eq:Pdifur}
for
$m=n-k$,
we obtain the system
\begin{equation*}
\begin{cases}
f(x)+g(x) = \dfrac{(x-x^2)}{k+1}, \\
g'(x) = (2x-1)\dfrac{2(n-k-1)}{k+1}.
\end{cases}
\end{equation*}
 Therefore,
$$f'(x) = (1-2x) \dfrac{2(n-k)}{k+1}.$$
 Hence
$$
B'_{n,k}(x) = (1-2x) \dfrac{2(n-k)}{k+1} (P_{n}^{(k-n+1)})^2+(1-2x)(P_{n}^{(k-n+1)})^2 =
(1-2x)\dfrac{2n-k+1}{k+1}(P_{n}^{(k-n+1)})^2.
$$
 Thus, the function
$B_{n,k}(x)$
does not decrease
on the closed interval
$[0,{1}/{2}]$
and
does not increase
on the closed interval
$[{1}/{2},1]$,
while the function
$P_{n}^{(k-n+1)}$
is zero
at the maximum points of
$A^2_{n,k}(x)$,
whence we obtain the assertion of the theorem.
\end{proof}

\end{document}
